\chapter{Moteur de recherche}

Une fois que nous avons des fichiers d'index inversées, et une pipline qui permet de transformer des requêtes naturelles en requêtes booléennes, 
il suffit de contruire le moteur de recherche qui trouve les documents pertinents a partir de la requête. Cette partie est prise en charge par la classe \texttt{SearchEngine} qui coordonne toute les étapes.

\section{Chargement de l'index}

La classe \texttt{SearchEngine} doit charger tout l'index en mémoire à son initialisation afin de pouvoir l'utiliser. Nous avons choisi de ranger les index comme un dictionnaire de \texttt{set} Python avec comme clé le token et comme valeur l'ensemble des ids des document contenant ce token.
Cela permet d'abord de rapidement accéder au tokens d'intéret.

